\documentclass[11pt]{article}
% Bien dich: pdflatex (hoac Overleaf). Moi hinh ve viet bang TikZ/pgfplots ngay trong file, khong can anh ngoai.
\usepackage[T5]{fontenc}
\usepackage[utf8]{inputenc}
\usepackage[vietnamese]{babel}
\usepackage[a4paper,margin=2.3cm]{geometry}
\usepackage{amsmath,amssymb}
\usepackage{booktabs}
\usepackage{array}
\usepackage{float}
\usepackage{caption}
\usepackage{xcolor}
\usepackage{tikz}
\usetikzlibrary{positioning,arrows.meta,shapes.geometric,fit,calc}
\usepackage{pgfplots}
\pgfplotsset{compat=1.17}
\usepackage[hidelinks]{hyperref}

\definecolor{cOrig}{HTML}{7F7F7F}
\definecolor{cV1}{HTML}{1F77B4}
\definecolor{cV2}{HTML}{D62728}
\definecolor{cV3}{HTML}{2CA02C}
\definecolor{cBox}{HTML}{EAF2FB}
\definecolor{cNew}{HTML}{FDECEA}
\definecolor{cNew3}{HTML}{E8F5E9}

\newcommand{\R}{\mathbb{R}}
\newcommand{\bmu}{\boldsymbol{\mu}}
\newcommand{\bSigma}{\boldsymbol{\Sigma}}
\newcommand{\vf}{\mathbf{f}}
\newcommand{\tr}{\operatorname{tr}}

\title{Báo cáo dự án khoa học\\[4pt]
\large Hai phiên bản tốt nhất hiện tại:\\
V2 với tổ hợp bộ phân loại, và V3-specgrid với một bộ phân loại CA}
\author{Thị Thanh Nga Hoàng}
\date{November 2026}

\begin{document}
\maketitle

\begin{abstract}
Báo cáo tổng hợp hai phiên bản cho kết quả cao nhất trong dự án học tăng dần theo lớp (CIL) không lưu mẫu cũ trên
ViT-B/16 với adapter. (A) \textbf{Bản tốt nhất có tổ hợp} là \textbf{V2}: backbone RSIAT bỏ $\mathcal{L}_{RS}$ và
$\mathcal{L}_{orth}$ (giữ RAE), bù drift bằng \textbf{STS} chạy một lần cuối mỗi task, và phân loại bằng tổ hợp z-score
của ba bộ phân loại CA, LDA, AGM. Trên ImageNet-A seed 1993, V2 đạt $75.63/66.69$ (trung bình/task cuối) so với
$74.76/65.77$ của RSIAT gốc; trung bình 3 seed $73.62/65.94$ so với $72.68/65.02$. (B) \textbf{Bản tốt nhất không tổ hợp}
(chỉ bộ phân loại CA) là \textbf{V3-specgrid}: giữ nguyên backbone và STS của V2, nhưng covariance lớp mới được
\emph{vay hình dạng} từ các lớp lân cận và cộng thêm \emph{độ bất định tích luỹ} do bù drift. Trên ImageNet-A seed 1993,
CA của V3-specgrid đạt $75.13/65.96$, cao hơn CA của V2 ($75.03/66.29$) ở trung bình nhưng thấp hơn ở task cuối. Chênh lệch
này nằm trong mức nhiễu của một seed; báo cáo nêu rõ điều đó.
\end{abstract}

\tableofcontents
\newpage

% =====================================================================================================================
\section{Tóm tắt kết quả}\label{sec:summary}

Mọi số trong báo cáo là độ chính xác top-1 (\%) trên tất cả các lớp đã thấy. ``Avg'' là trung bình qua các task, ``Last''
là sau task cuối.

\begin{table}[H]
\centering\small
\caption{Hai bản tốt nhất trên ImageNet-A, seed 1993 (Avg / Last). Cùng một backbone: hai bản cho đặc trưng test giống
hệt nhau ở cả 10 task; khác biệt chỉ nằm ở thống kê lớp và bộ phân loại.}
\label{tab:summary}
\begin{tabular}{llcc}
\toprule
Loại & Phiên bản & Avg / Last & So với RSIAT gốc \\
\midrule
Tham chiếu & RSIAT gốc (CA) & 74.76 / 65.77 & -- \\
\midrule
\textbf{Có tổ hợp} & \textbf{V2}, z(CA)+z(LDA)+z(AGM) & \textbf{75.63 / 66.69} & $+0.87$ / $+0.92$ \\
 & V3-specgrid, z(CA)+z(LDA)+z(AGM) & 75.54 / 65.77 & $+0.78$ / $\phantom{+}0.00$ \\
\midrule
\textbf{Không tổ hợp (CA)} & \textbf{V3-specgrid}, CA & \textbf{75.13} / 65.96 & $+0.37$ / $+0.19$ \\
 & V2, CA & 75.03 / \textbf{66.29} & $+0.27$ / $+0.52$ \\
 & V1, CA & 74.77 / 65.70 & $+0.01$ / $-0.07$ \\
\bottomrule
\end{tabular}
\end{table}

\paragraph{Đọc bảng.}
\begin{itemize}
\item \textbf{Có tổ hợp:} V2 là bản tốt nhất ở cả Avg và Last. V2 cũng là bản duy nhất đã chạy trên nhiều dataset và nhiều
      seed (mục~\ref{sec:res-v2}).
\item \textbf{Không tổ hợp:} V3-specgrid có Avg cao nhất ($75.13$, hơn V2 $+0.10$), nhưng V2 có Last cao hơn ($66.29$ so
      với $65.96$). Trên ImageNet-A, một ảnh test tương ứng $0.06$--$0.3$ điểm tuỳ task, nên $+0.10$ ở Avg chưa phải
      chênh lệch có ý nghĩa thống kê. V3-specgrid mới chạy một seed trên ImageNet-A.
\item Nếu cần \emph{một} bản để báo cáo chính, V2 có tổ hợp là lựa chọn an toàn nhất.
\end{itemize}

% =====================================================================================================================
\section{Bài toán và ký hiệu}\label{sec:setup}

Dữ liệu đến theo $T$ task; task $t$ mang tập lớp $\mathcal{C}_t$ rời nhau và tập ảnh $\mathcal{D}_t$. Khi học task
$t$ \emph{không} được giữ ảnh của các task trước (exemplar-free). Mô hình gồm backbone ViT-B/16 đông cứng với adapter
học được, cho đặc trưng $\vf(x)\in\R^d$ ($d=768$). Sau mỗi task, mô hình phải phân loại đúng mọi lớp đã thấy
$\mathcal{C}_{\le t}$.

Với mỗi lớp $c$, ta lưu \emph{thống kê lớp}: mean $\bmu_c$ và covariance $\bSigma_c$ của đặc trưng, ước lượng từ
$n_c$ ảnh train lúc lớp được học. Ảnh cũ không còn, nên khi adapter tiếp tục thay đổi ở task sau, đặc trưng ``trôi''
(\emph{drift}) và thống kê đã lưu không còn khớp. Mọi bộ phân loại trong báo cáo được dựng từ $(\bmu_c,\bSigma_c)$,
nên chất lượng của chúng phụ thuộc vào hai nguồn sai số:
\begin{enumerate}
\item[(I)] \textbf{biểu diễn}: các lớp có tách được trong không gian đặc trưng hiện tại hay không;
\item[(II)] \textbf{ước lượng thống kê}: $(\bmu_c,\bSigma_c)$ đã lưu có đúng với đặc trưng hiện tại hay không.
\end{enumerate}
V2 và V3 giữ nguyên backbone (phần I) và chỉ tác động lên phần (II) và cách dùng thống kê.

% =====================================================================================================================
\section{Pipeline chung của hai bản}\label{sec:common}

\begin{figure}[H]
\centering
\begin{tikzpicture}[
  box/.style={draw, rounded corners=2pt, fill=cBox, align=center, font=\small, minimum height=0.95cm, text width=2.5cm},
  newbox/.style={box, fill=cNew},
  v3box/.style={box, fill=cNew3},
  arr/.style={-{Stealth[length=2.2mm]}, thick}]
\node[box] (data) {Ảnh task $t$\\$\mathcal{D}_t$};
\node[box, right=0.55cm of data] (train) {Huấn luyện adapter\\CosFace + RAE};
\node[newbox, right=0.55cm of train] (sts) {STS cuối task:\\bám drift lớp cũ};
\node[box, right=0.55cm of sts] (stat) {Thống kê lớp mới\\$\bmu_c,\bSigma_c$};
\node[v3box, below=0.9cm of stat] (v3) {V3: vay hình dạng $S_c$\\+ bất định $U_c$};
\node[box, below=0.9cm of sts] (ca) {CA: huấn luyện lại\\head trên mẫu Gauss};
\node[newbox, below=0.9cm of train, text width=3.1cm] (ens) {V2: z(CA)+z(LDA)+z(AGM)\\V3-specgrid: chỉ CA};
\draw[arr] (data) -- (train);
\draw[arr] (train) -- (sts);
\draw[arr] (sts) -- (stat);
\draw[arr, dashed, cV3!70!black] (stat) -- (v3);
\draw[arr] (v3) -- (ca);
\draw[arr] (stat.south west) ++(0.1,0) |- ([yshift=0.25cm]ca.east);
\draw[arr] (ca) -- (ens);
\end{tikzpicture}
\caption{Pipeline cho một task $t\ge1$. Ô đỏ nhạt là phần V2 thêm vào so với RSIAT (STS thay SDC; tổ hợp ba bộ phân
loại). Ô xanh lá là phần V3-specgrid thêm vào so với V2 (thống kê lớp có độ bất định). Hai bản dùng chung mọi ô còn lại.}
\label{fig:pipeline}
\end{figure}

\subsection{Backbone và hàm mất mát}
Backbone giống RSIAT \cite{ssiat}: ViT-B/16 tiền huấn luyện ImageNet-21k, đông cứng, cộng adapter (FFN thắt cổ chai, 64
chiều) ở mỗi khối. Head là bộ phân loại cosine. Ở mỗi task, adapter được học với:
\begin{equation}
\mathcal{L} = \mathcal{L}_{\text{CosFace}} + \beta\,\mathcal{L}_{\text{RAE}},
\end{equation}
\begin{itemize}
\item $\mathcal{L}_{\text{CosFace}}$ \cite{cosface}: cross-entropy với lề cosine $m$, hệ số $s$, chỉ tính trên logit của
      lớp mới (lớp cũ không còn ảnh).
\item $\mathcal{L}_{\text{RAE}}$: bộ tự mã hoá dư $g(u)=u+\sigma(\mathrm{Dec}(\mathrm{Enc}(u)))$ học chung, ràng buộc đặc
      trưng mới vẫn tái tạo được đặc trưng của mô hình trước. Đây là cơ chế giữ duy nhất của backbone; bỏ RAE làm mất
      $3.6$ điểm ở task cuối trên ImageNet-A (đã kiểm chứng ở báo cáo V2).
\end{itemize}
So với RSIAT gốc, \textbf{cả hai bản đều bỏ} $\mathcal{L}_{RS}$ (lề quan hệ cặp, task 0) và $\mathcal{L}_{orth}$
($\lambda_{RS}=\gamma=0$). Hai thành phần này không cải thiện kết quả khi đã có các bước sau và làm pipeline phức tạp
hơn; mã V3 có \texttt{assert} để đảm bảo điều này.

\subsection{Thống kê lớp mới}
Sau khi học task $t$, với mỗi lớp mới $c\in\mathcal{C}_t$ có $n_c$ ảnh train ($n_c\approx30$ trên ImageNet-A):
\begin{equation}
\bmu_c=\frac1{n_c}\sum_i \vf(x_i),\qquad \hat\bSigma_c=\frac1{n_c-1}\sum_i(\vf(x_i)-\bmu_c)(\vf(x_i)-\bmu_c)^\top.
\end{equation}
V2 dùng $\bSigma_c=\hat\bSigma_c+10^{-3}I$. Vì $n_c\ll d=768$, $\hat\bSigma_c$ có hạng tối đa $n_c-1$: covariance suy biến
ở hơn $96\%$ số chiều. Đây là điểm V3 sửa (mục~\ref{sec:v3}).

\subsection{STS: bám drift của thống kê lớp cũ}\label{sec:sts}
RSIAT dùng SDC \cite{sdc}: dời mean lớp cũ theo trung bình có trọng số độ dịch của ảnh mới. SDC có hai điểm yếu: trọng
số không phân biệt ảnh mới nào thực sự gần lớp cũ, và covariance lớp cũ không được cập nhật. STS (Self-Tracking
Statistics) thay SDC, chạy \textbf{một lần cuối mỗi task}, với các bước sau.

\paragraph{Bước 1: tập dò.} Đưa ảnh của task $t$ qua mô hình trước ($\vf_{t-1}$) và mô hình hiện tại ($\vf_t$), thu
cặp $(\vf_{t-1}(x),\vf_t(x))$. Với lớp cũ có quá ít ảnh liên quan (khối lượng $N_c<5$), sinh thêm $16$ mẫu đối kháng
theo ADC \cite{adc}: đẩy ảnh mới về phía $\bmu_c$ trong không gian cũ bằng $10$ bước gradient, bước $2/255$.

\paragraph{Bước 2: bước E có thành phần nền.} Ảnh mới thường không thuộc lớp cũ nào. Dùng mô hình trộn gồm các Gauss
lớp cũ $\mathcal{N}(\bmu_c,\bSigma_c)$ và một thành phần nền có xác suất tiên nghiệm $\pi_0=0.5$:
\begin{equation}
r_{ic}=\frac{\pi_c\,\mathcal{N}(\vf_{t-1}(x_i);\bmu_c,\bSigma_c)^{\eta}}
{\pi_0\,p_0(x_i)+\sum_{c'}\pi_{c'}\,\mathcal{N}(\vf_{t-1}(x_i);\bmu_{c'},\bSigma_{c'})^{\eta}},\qquad \eta=0.3.
\end{equation}
Số mũ $\eta<1$ làm mềm likelihood trong không gian $768$ chiều. Ảnh không gần lớp cũ nào được nền ``hấp thụ'' thay vì bị
ép gán cho lớp gần nhất. $N_c=\sum_i r_{ic}$ là khối lượng ảnh liên quan tới lớp $c$.

\paragraph{Bước 3: ánh xạ drift.} Khớp ánh xạ toàn cục $\vf_t\approx \vf_{t-1}(I+B)+b$ bằng Procrustes \cite{procrustes}
(phần quay trực giao) trên các cặp có trọng số, với độ co $\alpha\in\{10^{-2},\dots,10^{2}\}$ chọn bằng split-half: khớp
trên một nửa tập dò, đo sai số dự đoán trên nửa kia. Với \texttt{procrustes}, mọi lớp dùng chung ánh xạ toàn cục; dịch
chuyển riêng từng lớp co về ánh xạ chung với $n_0=50$.

\paragraph{Bước 4: vận chuyển.} Mean và covariance lớp cũ đi qua cùng ánh xạ $A=I+B$:
\begin{equation}
\bmu_c\leftarrow A^\top\bmu_c+b,\qquad \bSigma_c\leftarrow A^\top\bSigma_c A.
\end{equation}
Khác SDC, covariance cũng được cập nhật. Đây là lý do LDA và AGM (hai bộ phân loại dùng covariance) tăng nhiều nhất từ V1
lên V2 ($+0.66$ và $+0.85$ ở task cuối).

\subsection{Bộ phân loại CA}
Theo SLCA \cite{slca}, head cosine được huấn luyện lại $10$ epoch trên mẫu Gauss sinh từ thống kê của mọi lớp đã thấy:
$\tilde\vf\sim\mathcal{N}(\kappa_c\bmu_c,\bSigma^{\text{head}}_c)$, $256$ mẫu mỗi lớp mỗi epoch. Hệ số co mean
$\kappa_c=0.9+\delta_c$ với $\delta_c=0.1\,(t_c+1)/(t+1)$, trong đó $t_c$ là task học lớp $c$: lớp càng cũ, mean càng
bị co về gốc, bù cho việc thống kê của lớp cũ kém tin cậy hơn. V2 dùng $\bSigma^{\text{head}}_c=\bSigma_c$; V3-specgrid
dùng covariance dự đoán $\bSigma^{\text{pred}}_c$ (mục~\ref{sec:v3-pred}).

% =====================================================================================================================
\section{Bản A -- V2 với tổ hợp ba bộ phân loại}\label{sec:v2}

\subsection{Động cơ}
Ba bộ phân loại dựng từ cùng $(\bmu_c,\bSigma_c)$ nhưng dùng chúng theo cách khác nhau, nên sai số của chúng không
trùng nhau. Tổ hợp tận dụng điều này mà không tốn thêm huấn luyện backbone.

\subsection{Ba bộ phân loại}
\paragraph{CA.} Như mục trên; logit $\ell^{CA}_c(x)=\cos(\vf(x),w_c)$.

\paragraph{LDA.} Phân tích biệt thức tuyến tính \cite{esl} với covariance gộp, chính quy hoá ridge:
\begin{equation}
\bar\bSigma=\frac1{|\mathcal{C}_{\le t}|}\sum_c\bSigma_c,\quad
\bSigma_\lambda=\bar\bSigma+\lambda I,\ \lambda=\tfrac1d\tr\bar\bSigma,\quad
\ell^{LDA}_c(x)=\bmu_c^\top\bSigma_\lambda^{-1}\vf(x)-\tfrac12\bmu_c^\top\bSigma_\lambda^{-1}\bmu_c.
\end{equation}
LDA là bộ phân loại Bayes khi các lớp là Gauss cùng covariance; nó không cần huấn luyện và chỉ phụ thuộc vào thống kê.

\paragraph{AGM (Analytic Gaussian-Moment head).} Hồi quy ridge trên đặc trưng ngẫu nhiên \cite{rahimi}
$\phi(\vf)=[\vf;\ \mathrm{ReLU}(W\vf);\ 10]\in\R^{M}$, $M=4769$, $W$ ngẫu nhiên cố định. Thay vì sinh mẫu, các mô-men cần
cho nghiệm đóng được tính \emph{chính xác} từ Gauss của mỗi lớp (kỳ vọng của ReLU của biến Gauss có dạng đóng):
\begin{equation}
G=\sum_c n_c\,\mathbb{E}_{\vf\sim\mathcal{N}(\bmu_c,\bSigma_c)}\!\left[\phi\phi^\top\right],\quad
H_{:,c}=n_c\,\mathbb{E}[\phi],\quad
\Omega=(G+\lambda I)^{-1}H,\quad \lambda=0.1\,\tfrac{\tr G}{M}.
\end{equation}
Logit $\ell^{AGM}(x)=\phi(\vf(x))^\top\Omega$. AGM là bộ phân loại phi tuyến, không có nhiễu lấy mẫu như CA.

\subsection{Tổ hợp z-score}
Ba loại logit có thang đo khác nhau (cosine, log-mật độ, hồi quy), nên trước khi cộng, mỗi logit được chuẩn hoá theo
từng ảnh:
\begin{equation}
z(\ell)_c=\frac{\ell_c-\mathrm{mean}_{c'}\ell_{c'}}{\mathrm{std}_{c'}\ell_{c'}},\qquad
\hat y(x)=\arg\max_c\left[z(\ell^{CA})_c+z(\ell^{LDA})_c+z(\ell^{AGM})_c\right].
\end{equation}
Trọng số bằng nhau, không có tham số nào chọn trên tập test.

\subsection{Cấu hình}
\texttt{exps/sts/imageneta\_raests\_end\_s1993.json} (ImageNet-A); \texttt{exps/v2/*} cho các dataset khác. Các khoá
chính: \texttt{rae: true, sts: true, sts\_freq: end, sts\_map: procrustes, lambda\_rs: 0, gamma: 0}. Kết quả tổ hợp tính
bằng \texttt{exps/agm/final\_eval.py} trên dump của từng task.

% =====================================================================================================================
\section{Bản B -- V3-specgrid: thống kê lớp có độ bất định}\label{sec:v3}

\subsection{Động cơ}
V2 có hai điểm yếu về thống kê:
\begin{enumerate}
\item \textbf{Covariance lớp mới suy biến.} Với $n_c\approx30$ ảnh trong $768$ chiều, $\hat\bSigma_c$ chỉ có $\approx29$
      hướng khác $0$; ridge $10^{-3}I$ là một sửa tạm. Mẫu Gauss dùng cho CA vì vậy chỉ trải trên một không gian con nhỏ.
\item \textbf{Bù drift không có sai số.} STS cho một ước lượng điểm của $\bmu_c$ sau drift, nhưng ước lượng này có sai
      số và sai số tích luỹ qua các task. V2 coi thống kê lớp cũ là chính xác.
\end{enumerate}
V3 tách covariance của mỗi lớp thành hai phần có vai trò khác nhau:
\begin{equation}
\bSigma^{\text{pred}}_c=\underbrace{\Big(1+\tfrac1{n_c}\Big)S_c}_{\text{độ phân tán của lớp}}+\underbrace{U_c}_{\text{sai số của ước lượng}},
\end{equation}
trong đó $S_c$ là \emph{hình dạng} của lớp (phân tán của đặc trưng quanh mean), $U_c$ là \emph{độ bất định} của mean
đã lưu. Hệ số $1+1/n_c$ là phần bất định của mean ước lượng từ $n_c$ mẫu. Đây là dạng phân bố dự đoán của mô hình Gauss
có tham số ước lượng \cite{murphy}.

\subsection{Vay hình dạng từ các lớp lân cận}\label{sec:v3-borrow}
Các lớp gần nhau trong không gian đặc trưng thường có hình dạng phân tán tương tự. V3 dùng tiên nghiệm cho $S_c$ là
trung bình có trọng số các covariance nguồn (mọi lớp cũ và các lớp mới khác của cùng task, trừ chính $c$):
\begin{equation}
P_c=\sum_{j\ne c}w_{cj}\,S_j,\qquad w_{cj}=\frac{\exp\big(\tau\cos(\bmu_c,\bmu_j)\big)}{\sum_{j'\ne c}\exp\big(\tau\cos(\bmu_c,\bmu_{j'})\big)}.
\end{equation}
Ước lượng hậu nghiệm co $\hat\bSigma_c$ về $P_c$ với độ mạnh $n_0$ ``ảnh giả'' (dạng liên hợp Wishart nghịch đảo
\cite{murphy}):
\begin{equation}
S_c=\frac{n_0\,P_c+(n_c-1)\,\hat\bSigma_c}{n_0+n_c-1}+10^{-3}I.
\end{equation}
Với $n_0=0$, $S_c$ trùng với covariance của V2. $\tau$ điều khiển mức tập trung vào các lớp gần nhất ($\tau=0$: trung
bình đều).

\paragraph{Chọn $(n_0,\tau)$ bằng dữ liệu train.} Mỗi task, ảnh của lớp mới (lớp có ít nhất $4$ ảnh) được chia đôi ngẫu
nhiên. Với mỗi cặp $(n_0,\tau)$ trên lưới, ước lượng $(\bmu_c,S_c)$ từ một nửa và tính log-likelihood Gauss của nửa kia,
rồi đổi vai (2-fold). Chọn cặp có log-likelihood trung bình mỗi ảnh lớn nhất. Không dùng tập test. Lưới
\textbf{specgrid}:
\[
n_0\in\{0,5,10,20,50,100\},\qquad \tau\in\{4,8,16,32\}.
\]
\textbf{Kết quả chọn:} $n_0=100$, $\tau=4$ ở cả 10 task. Log-likelihood giữ lại tăng rất mạnh so với $n_0=0$ (ví dụ task
1: $-2469.5$ so với $-776717$ mỗi ảnh): covariance của V2 gần như suy biến nên gán mật độ cực thấp cho ảnh ngoài mẫu.
Lưu ý $n_0=100$ nằm ở biên lưới (xem mục~\ref{sec:limits}).

\subsection{Độ bất định tích luỹ $U_c$}\label{sec:v3-u}
Lớp mới có $U_c=0$. Mỗi lần STS vận chuyển lớp $c$ bằng ánh xạ $A_c$, hình dạng và độ bất định cập nhật theo bước dự
đoán của bộ lọc Kalman:
\begin{equation}
S_c\leftarrow A_c^\top S_cA_c,\qquad U_c\leftarrow A_c^\top U_cA_c+Q_c.
\end{equation}
$Q_c$ là sai số của chính ánh xạ drift, đo \emph{ngoài mẫu} bằng cross-fitting: chia tập dò thành hai fold, khớp ánh xạ
trên fold này và đo phần dư $\vf_t(x)-\mathrm{pred}(x)$ trên fold kia. Từ phần dư có trọng số $r_{ic}$:
\begin{equation}
R_c=\text{covariance phần dư ngoài mẫu},\qquad Q_c=R_c\Big(1+\tfrac1{n_{\text{eff}}}\Big),
\end{equation}
với $n_{\text{eff}}$ là số ảnh hiệu dụng. Lớp không có ánh xạ riêng dùng phần dư toàn cục $R_g$ với $n_{\text{eff}}$ là
số ảnh thật. Để ổn định trong $768$ chiều, $Q$ được xấp xỉ bằng PPCA: giữ $16$ hướng chính cộng một phần đẳng hướng,
bảo toàn vết.

\subsection{Dùng $\bSigma^{\text{pred}}$ ở đâu}\label{sec:v3-pred}
Trong V3-specgrid, $\bSigma^{\text{pred}}_c=(1+1/n_c)S_c+U_c$ thay cho $\bSigma_c$ ở \emph{mọi} chỗ dùng covariance:
\begin{itemize}
\item bước E của STS (mục~\ref{sec:sts}): lớp cũ có bất định lớn có mật độ trải rộng hơn;
\item sinh mẫu Gauss cho CA: lớp cũ có $U_c$ lớn được biểu diễn bằng đám mẫu rộng hơn, head buộc phải chừa lề cho sai số
      của mean;
\item LDA và AGM (khi dùng tổ hợp).
\end{itemize}
Mean, ánh xạ drift và backbone giống hệt V2. Đã kiểm chứng bằng \texttt{exps/v3/check\_backbone.py}: đặc trưng test của
V3-specgrid và V2 trùng nhau từng bit ở cả 10 task.

\subsection{Cấu hình}
\texttt{exps/v3/imageneta\_v3specgrid\_s1993.json}: như V2 cộng \texttt{v3: true}, \texttt{v3\_n0\_grid},
\texttt{v3\_tau\_grid}, \texttt{v3\_select: loglik}, \texttt{v3\_estep: pred}, \texttt{v3\_head\_cov: pred}. Ba khoá cuối
là hành vi của mã lúc chạy V3-specgrid; mã hiện tại có mặc định khác (của V3.1), nên đã ghi rõ chúng vào cấu hình để
chạy lại cho đúng.

% =====================================================================================================================
\section{Thiết lập thực nghiệm}\label{sec:exp}

\textbf{ImageNet-A}: 200 lớp, 10 task $\times$ 20 lớp, $\approx30$ ảnh train/lớp, thứ tự lớp xáo theo seed.
\textbf{VTAB}: 50 lớp, 5 task (mỗi task một miền). \textbf{CUB-200}: 200 lớp, 20 task $\times$ 10 lớp. Siêu tham số
giống RSIAT gốc (ImageNet-A: 20/15 epoch, CA 10 epoch, batch 48, $s=30$, $m=0.3$, RAE $\beta=3$, \texttt{code\_dims}
128). STS: $\eta=0.3$, $\pi_0=0.5$, ngưỡng $N_c=5$, $n_0^{STS}=50$, Procrustes.

\paragraph{Độ tin cậy.} Huấn luyện tất định. Tập test ImageNet-A nhỏ: một ảnh bằng $0.06$--$0.3$ điểm, nên chênh lệch
dưới $\approx0.3$ trên một seed nằm trong mức nhiễu. RSIAT gốc trên CUB là số do tác giả cung cấp. Số CA trong bảng so
sánh không tổ hợp lấy từ log huấn luyện; số tổ hợp tính offline trên dump (CA tính lại từ dump có thể lệch log
$\le0.1$ điểm, ví dụ CA của V2: $75.03/66.29$ trong log, $75.02/66.23$ trên dump).

% =====================================================================================================================
\section{Kết quả bản A (V2, có tổ hợp)}\label{sec:res-v2}

\subsection{ImageNet-A}

\begin{table}[H]
\centering\small
\caption{ImageNet-A, seed 1993, có tổ hợp. RSIAT gốc: CA; các bản còn lại: z(CA)+z(LDA)+z(AGM). In đậm: cao nhất mỗi cột.}
\label{tab:ens1993}
\setlength{\tabcolsep}{3.6pt}
\begin{tabular}{l*{10}{c}cc}
\toprule
Task & 0 & 1 & 2 & 3 & 4 & 5 & 6 & 7 & 8 & 9 & Avg & Last \\
\midrule
RSIAT gốc & 86.29 & 82.50 & 79.62 & \textbf{78.65} & 75.51 & 73.23 & 70.56 & 68.59 & 66.88 & 65.77 & 74.76 & 65.77 \\
V1 & 87.43 & 84.72 & \textbf{81.09} & 79.08 & 75.03 & 74.36 & 71.01 & 68.51 & 66.81 & 65.90 & 75.39 & 65.90 \\
\textbf{V2} & 87.43 & 84.44 & 80.88 & 78.51 & \textbf{75.62} & \textbf{74.87} & \textbf{71.63} & \textbf{68.99} & \textbf{67.24} & \textbf{66.69} & \textbf{75.63} & \textbf{66.69} \\
V3-specgrid & \textbf{88.57} & \textbf{85.56} & 80.25 & \textbf{79.08} & 75.39 & 74.15 & 71.10 & 68.91 & 66.67 & 65.77 & 75.54 & 65.77 \\
\bottomrule
\end{tabular}
\end{table}
(Task 3: V1 và V3-specgrid cùng $79.08$, cao nhất.)

\begin{figure}[H]
\centering
\begin{tikzpicture}
\begin{axis}[width=0.92\linewidth, height=6.2cm, xlabel={Task}, ylabel={Độ chính xác (\%)},
  xtick={0,...,9}, ymin=64, ymax=89.5, grid=major, grid style={black!12},
  legend style={at={(0.02,0.03)}, anchor=south west, font=\small, fill=white, fill opacity=0.9}]
\addplot[cOrig, thick, mark=o] coordinates {(0,86.29)(1,82.50)(2,79.62)(3,78.65)(4,75.51)(5,73.23)(6,70.56)(7,68.59)(8,66.88)(9,65.77)};
\addlegendentry{RSIAT gốc (CA)}
\addplot[cV1, thick, mark=square] coordinates {(0,87.43)(1,84.72)(2,81.09)(3,79.08)(4,75.03)(5,74.36)(6,71.01)(7,68.51)(8,66.81)(9,65.90)};
\addlegendentry{V1 (tổ hợp)}
\addplot[cV3, thick, dashed, mark=triangle] coordinates {(0,88.57)(1,85.56)(2,80.25)(3,79.08)(4,75.39)(5,74.15)(6,71.10)(7,68.91)(8,66.67)(9,65.77)};
\addlegendentry{V3-specgrid (tổ hợp)}
\addplot[cV2, very thick, mark=*] coordinates {(0,87.43)(1,84.44)(2,80.88)(3,78.51)(4,75.62)(5,74.87)(6,71.63)(7,68.99)(8,67.24)(9,66.69)};
\addlegendentry{V2 (tổ hợp)}
\end{axis}
\end{tikzpicture}
\caption{ImageNet-A seed 1993, có tổ hợp. V3-specgrid dẫn ở các task đầu (thống kê lớp mới tốt hơn) nhưng tụt dưới V2 từ
task 4 (độ bất định tích luỹ của lớp cũ bị ước lượng quá lớn, mục~\ref{sec:why}). V2 cao nhất từ task 4 tới cuối.}
\label{fig:ens}
\end{figure}

\begin{table}[H]
\centering\small
\caption{ImageNet-A seed 1993: từng bộ phân loại (Avg / Last, tính trên dump).}
\label{tab:heads}
\begin{tabular}{lccccc}
\toprule
 & CA & LDA & AGM & z(CA+AGM) & z(CA+LDA+AGM) \\
\midrule
V1 & 74.76 / 65.64 & 74.33 / 64.98 & 74.19 / 64.65 & 75.25 / 65.90 & 75.39 / 65.90 \\
V2 & 75.02 / \textbf{66.23} & 74.42 / \textbf{65.64} & \textbf{74.55 / 65.50} & \textbf{75.50 / 66.29} & \textbf{75.63 / 66.69} \\
V3-specgrid & \textbf{75.13} / 65.96 & \textbf{74.64} / 65.17 & 73.95 / 64.52 & 75.42 / 65.57 & 75.54 / 65.77 \\
\bottomrule
\end{tabular}
\end{table}
Tổ hợp luôn cao hơn bộ phân loại đơn tốt nhất ($+0.5$ Avg với V2). Ở V3-specgrid, AGM yếu đi rõ ($-0.6/-1.0$ so với V2),
kéo tổ hợp xuống.

\begin{table}[H]
\centering\small
\caption{ImageNet-A, 3 seed, có tổ hợp (Avg / Last). V2 thắng RSIAT gốc ở cả 3 seed.}
\label{tab:ina3}
\begin{tabular}{lcccc}
\toprule
Seed & RSIAT gốc & V1 & V2 & V2 $-$ gốc \\
\midrule
1993 & 74.76 / 65.77 & 75.39 / 65.90 & \textbf{75.63 / 66.69} & $+0.87$ / $+0.92$ \\
1994 & 72.18 / 63.99 & 72.98 / 64.98 & \textbf{73.06 / 65.04} & $+0.88$ / $+1.05$ \\
1995 & 71.10 / 65.31 & \textbf{72.19} / 65.70 & 72.18 / \textbf{66.10} & $+1.08$ / $+0.79$ \\
\midrule
\textbf{Trung bình} & 72.68 / 65.02 & 73.52 / 65.53 & \textbf{73.62 / 65.94} & $\mathbf{+0.94}$ / $\mathbf{+0.92}$ \\
\bottomrule
\end{tabular}
\end{table}

\begin{figure}[H]
\centering
\begin{tikzpicture}
\begin{axis}[ybar, bar width=12pt, width=0.8\linewidth, height=5.6cm, ymin=64, ymax=74.5,
  symbolic x coords={Avg,Last}, xtick=data, enlarge x limits=0.5, ylabel={Độ chính xác (\%)},
  nodes near coords, nodes near coords style={font=\scriptsize, rotate=90, anchor=west},
  legend style={at={(0.5,1.03)}, anchor=south, legend columns=3, font=\small}]
\addplot[fill=cOrig!70] coordinates {(Avg,72.68)(Last,65.02)};
\addplot[fill=cV1!70] coordinates {(Avg,73.52)(Last,65.53)};
\addplot[fill=cV2!80] coordinates {(Avg,73.62)(Last,65.94)};
\legend{RSIAT gốc, V1, V2}
\end{axis}
\end{tikzpicture}
\caption{ImageNet-A, trung bình 3 seed (có tổ hợp).}
\end{figure}

\subsection{VTAB}
\begin{table}[H]
\centering\small
\caption{VTAB (một lần chạy tất định; seed không ảnh hưởng vì thứ tự miền cố định). RSIAT gốc: CA; V1, V2: tổ hợp.}
\begin{tabular}{l*{5}{c}cc}
\toprule
Task & 0 & 1 & 2 & 3 & 4 & Avg & Last \\
\midrule
RSIAT gốc & 96.70 & 93.93 & 94.00 & \textbf{94.27} & 94.69 & 94.72 & 94.69 \\
V1 & \textbf{97.10} & 94.29 & 94.17 & 94.14 & 94.86 & 94.91 & 94.86 \\
\textbf{V2} & \textbf{97.10} & \textbf{94.36} & \textbf{94.34} & 94.25 & \textbf{94.91} & \textbf{94.99} & \textbf{94.91} \\
\bottomrule
\end{tabular}
\end{table}
Từng bộ phân loại của V2 trên VTAB (Avg / Last): CA $94.80/94.73$, LDA $93.99/93.94$, AGM $94.34/94.42$, tổ hợp
$94.99/94.91$.

\subsection{CUB-200}
\begin{table}[H]
\centering\small
\caption{CUB, seed 1993 (Avg / Last). RSIAT gốc do tác giả cung cấp.}
\label{tab:cub}
\begin{tabular}{lccccc}
\toprule
 & CA & LDA & AGM & z(CA+AGM) & z(CA+LDA+AGM) \\
\midrule
RSIAT gốc & 93.93 / 90.37 & -- & -- & -- & -- \\
V1 & 93.99 / 90.46 & 93.16 / 88.72 & 93.40 / 89.06 & 93.63 / 89.40 & 93.66 / 89.57 \\
V2 & \textbf{94.16 / 90.54} & 93.38 / 89.10 & 93.50 / 89.65 & 93.67 / 89.99 & 93.69 / 90.08 \\
\bottomrule
\end{tabular}
\end{table}
Trên CUB (phân loại chim rất mịn), LDA và AGM yếu hơn CA $1$--$1.5$ điểm ở task cuối, nên tổ hợp trọng số bằng nhau thấp
hơn CA. Trên CUB, \textbf{bản tốt nhất là CA của V2} ($94.16/90.54$, vượt RSIAT gốc $+0.23/+0.17$); tổ hợp không vượt
RSIAT gốc ở task cuối.

\begin{figure}[H]
\centering
\begin{tikzpicture}
\begin{axis}[width=0.92\linewidth, height=6cm, xlabel={Task}, ylabel={Độ chính xác (\%)},
  xtick={0,2,...,19}, ymin=89.3, ymax=100.3, grid=major, grid style={black!12},
  legend style={at={(0.98,0.97)}, anchor=north east, font=\small}]
\addplot[cOrig, thick, mark=o, mark size=1.5pt] coordinates {(0,100.00)(1,98.79)(2,98.55)(3,98.02)(4,97.50)(5,95.24)(6,93.70)(7,94.23)(8,93.82)(9,93.27)(10,93.40)(11,92.58)(12,91.79)(13,91.65)(14,91.23)(15,91.43)(16,91.21)(17,91.02)(18,90.80)(19,90.37)};
\addlegendentry{RSIAT gốc (CA)}
\addplot[cV2, very thick, mark=*, mark size=1.5pt] coordinates {(0,100.00)(1,98.79)(2,98.55)(3,98.02)(4,97.68)(5,95.39)(6,94.09)(7,94.34)(8,94.20)(9,93.60)(10,93.40)(11,92.94)(12,92.19)(13,92.20)(14,91.68)(15,91.80)(16,91.36)(17,91.35)(18,91.02)(19,90.54)};
\addlegendentry{V2 (CA)}
\addplot[cV2, thick, dashed, mark=x, mark size=2pt] coordinates {(0,99.15)(1,98.38)(2,98.26)(3,97.80)(4,97.68)(5,95.39)(6,93.83)(7,94.12)(8,93.92)(9,93.27)(10,93.25)(11,92.37)(12,91.33)(13,91.29)(14,90.95)(15,90.96)(16,90.77)(17,90.70)(18,90.35)(19,90.08)};
\addlegendentry{V2 (tổ hợp)}
\end{axis}
\end{tikzpicture}
\caption{CUB, seed 1993. CA của V2 bằng hoặc cao hơn RSIAT gốc ở cả 20 task; tổ hợp (nét đứt) thấp hơn CA.}
\end{figure}

% =====================================================================================================================
\section{Kết quả bản B (V3-specgrid, không tổ hợp)}\label{sec:res-v3}

\begin{table}[H]
\centering\small
\caption{ImageNet-A, seed 1993, chỉ bộ phân loại CA (số lấy từ log huấn luyện). In đậm: cao nhất mỗi cột.}
\label{tab:ca1993}
\setlength{\tabcolsep}{3.6pt}
\begin{tabular}{l*{10}{c}cc}
\toprule
Task & 0 & 1 & 2 & 3 & 4 & 5 & 6 & 7 & 8 & 9 & Avg & Last \\
\midrule
RSIAT gốc & \textbf{86.29} & 82.50 & 79.62 & \textbf{78.65} & 75.51 & 73.23 & 70.56 & 68.59 & 66.88 & 65.77 & 74.76 & 65.77 \\
V1 & \textbf{86.29} & 82.22 & 79.62 & 78.37 & 75.15 & 73.74 & 70.65 & 68.91 & 67.03 & 65.70 & 74.77 & 65.70 \\
V2 & \textbf{86.29} & 82.22 & 79.62 & 78.51 & \textbf{75.62} & \textbf{74.05} & \textbf{71.19} & \textbf{69.07} & \textbf{67.46} & \textbf{66.29} & 75.03 & \textbf{66.29} \\
\textbf{V3-specgrid} & \textbf{86.29} & \textbf{84.17} & \textbf{80.46} & 78.37 & \textbf{75.62} & 73.64 & 70.74 & 68.83 & 67.17 & 65.96 & \textbf{75.13} & 65.96 \\
\midrule
V3 $-$ V2 & 0.00 & $+1.95$ & $+0.84$ & $-0.14$ & 0.00 & $-0.41$ & $-0.45$ & $-0.24$ & $-0.29$ & $-0.33$ & $+0.10$ & $-0.33$ \\
\bottomrule
\end{tabular}
\end{table}

\begin{figure}[H]
\centering
\begin{tikzpicture}
\begin{axis}[width=0.92\linewidth, height=5.6cm, xlabel={Task}, ylabel={V3-specgrid $-$ V2 (CA, điểm \%)},
  xtick={0,...,9}, ymin=-0.8, ymax=2.2, grid=major, grid style={black!12}, ybar, bar width=12pt]
\addplot[fill=cV3!70, draw=cV3!60!black] coordinates {(0,0)(1,1.95)(2,0.84)(3,-0.14)(4,0)(5,-0.41)(6,-0.45)(7,-0.24)(8,-0.29)(9,-0.33)};
\end{axis}
\end{tikzpicture}
\caption{Chênh lệch CA giữa V3-specgrid và V2 theo task. Lợi ích tập trung ở task 1--2; từ task 5, V3 thua đều
$0.2$--$0.5$ điểm.}
\label{fig:cadiff}
\end{figure}

\paragraph{Lựa chọn của tiêu chí.} Ở cả 10 task, tiêu chí 2-fold chọn $n_0=100$, $\tau=4$. Ở task 0, CA không phụ thuộc
covariance (head học trực tiếp trên ảnh thật), nên ba bản V1, V2, V3 bằng nhau ($86.29$).

\subsection{Vì sao V3-specgrid tăng ở đầu và giảm ở cuối}\label{sec:why}
\begin{itemize}
\item \textbf{Task 1--2 tăng ($+1.95$, $+0.84$).} Lúc này phần lớn lớp còn ``trẻ'', $U_c$ còn nhỏ, và lợi ích chính là
      covariance vay ($n_0=100$) không còn suy biến. Mẫu Gauss cho CA trải đúng hơn theo hình dạng thật của lớp, nên head
      phân tách lớp mới và lớp cũ tốt hơn.
\item \textbf{Từ task 5 giảm.} $U_c$ tích luỹ theo số lần vận chuyển. Chẩn đoán trên dump cho thấy $\tr U_c$ lớn gấp
      $4$--$7$ lần sai số bình phương thật của mean lớp cũ (đo được nhờ thống kê oracle), và sau khi làm trắng theo
      $S_c$, $U_c$ của các lớp cũ nhất lớn gấp $13$--$16$ lần so với lớp mới. Độ bất định bị ước lượng quá lớn làm đám
      mẫu của lớp cũ quá rộng; head CA chừa lề quá nhiều, độ chính xác lớp cũ giảm. Trong bước E, lớp cũ cũng nhận ít
      khối lượng hơn.
\item Hai hiệu ứng gần như triệt tiêu nhau: $+0.10$ ở Avg, $-0.33$ ở Last.
\end{itemize}

\subsection{Các biến thể V3 đã thử sau specgrid (ImageNet-A seed 1993)}
\begin{table}[H]
\centering\small
\caption{Các biến thể V3 (Avg / Last). Không biến thể nào vượt V2 ở tổ hợp; CA của specgrid là cao nhất ở Avg.}
\begin{tabular}{lcc}
\toprule
Biến thể & CA & Tổ hợp \\
\midrule
V2 (tham chiếu) & 75.02 / 66.23 & \textbf{75.63 / 66.69} \\
V3, lưới rộng ($n_0$ tới 2000, $\tau$ từ 0) & -- & 75.48 / 65.64 \\
\textbf{V3-specgrid} & \textbf{75.13} / 65.96 & 75.54 / 65.77 \\
V3.1: chọn bằng log-loss có điều kiện, bước E và head dùng $S$ & 75.10 / 65.90 & 75.11 / 65.17 \\
V3.2: thêm co RDA $\gamma$ & \multicolumn{2}{c}{dừng giữa chừng, không có kết quả} \\
\bottomrule
\end{tabular}
\end{table}
(Cột CA ở bảng này tính trên dump để so cùng điều kiện.) V3.1 chọn $n_0=0$ ở task 0 rồi $n_0=500$, $\tau=0$ ở task 1--9:
tiêu chí log-loss có điều kiện bị chi phối bởi vài lớp và đẩy lựa chọn ra biên. V3.2 bị dừng theo yêu cầu.

% =====================================================================================================================
\section{Thảo luận}\label{sec:discussion}

\paragraph{Nên dùng bản nào.}
\begin{itemize}
\item \textbf{Có tổ hợp: V2.} Cao nhất trên ImageNet-A (1 seed và 3 seed) và VTAB; ổn định qua seed ($+0.94/+0.92$ trung
      bình 3 seed so với RSIAT gốc).
\item \textbf{Không tổ hợp:} trên ImageNet-A, V3-specgrid và V2 ngang nhau trong mức nhiễu (V3 hơn $0.10$ ở Avg, V2 hơn
      $0.33$ ở Last). Nếu ưu tiên Last hoặc cần kết quả trên nhiều dataset, CA của V2 tốt hơn (trên CUB, CA của V2 là bản
      tốt nhất). V3-specgrid chỉ nên gọi là ``tốt nhất không tổ hợp'' theo tiêu chí Avg trên ImageNet-A seed 1993.
\end{itemize}

\paragraph{Giới hạn.}\label{sec:limits}
\begin{itemize}
\item V3-specgrid mới có một lần chạy (ImageNet-A, seed 1993); chưa chạy VTAB, CUB và seed khác.
\item $n_0=100$ được chọn ở biên của lưới specgrid; mở rộng lưới (V3, V3.1) chọn $n_0$ lớn hơn nhưng kết quả không tốt
      hơn. Tiêu chí log-likelihood giữ lại không phản ánh đúng độ chính xác phân loại.
\item Ước lượng $U_c$ quá lớn là nguyên nhân chính làm V3 giảm ở các task cuối. Cần hiệu chỉnh $Q_c$ (ví dụ đối chiếu với
      sai số thật trên các lớp có đủ ảnh dò) trước khi V3 có thể vượt V2.
\item Mức cải thiện so với RSIAT gốc nhỏ ($<1$ điểm). Trên CUB, tổ hợp trọng số bằng nhau làm giảm điểm.
\end{itemize}

\paragraph{Vị trí so với công trình liên quan.} STS kế thừa SDC \cite{sdc} (ước lượng drift từ ảnh mới), ADC \cite{adc}
(mẫu đối kháng) và LDC \cite{ldc} (ánh xạ đặc trưng cũ--mới); điểm riêng là bước E có thành phần nền và vận chuyển cả
covariance. V3 thêm hai ý: vay hình dạng covariance từ lớp lân cận (tương tự co Bayes) và theo dõi độ bất định của thống
kê như bộ lọc Kalman. Cần so trực tiếp với ADC và LDC trên cùng backbone để khẳng định đóng góp.

% =====================================================================================================================
\section{Kết luận}
\begin{itemize}
\item \textbf{Bản tốt nhất có tổ hợp: V2} (RSIAT bỏ $\mathcal{L}_{RS}$, $\mathcal{L}_{orth}$; giữ RAE; STS cuối task;
      z(CA)+z(LDA)+z(AGM)): ImageNet-A $75.63/66.69$ (seed 1993), $73.62/65.94$ (3 seed); VTAB $94.99/94.91$.
\item \textbf{Bản tốt nhất không tổ hợp (theo Avg): V3-specgrid} (V2 + covariance vay hình dạng + độ bất định tích luỹ,
      chỉ CA): ImageNet-A seed 1993 $75.13/65.96$, so với CA của V2 $75.03/66.29$ và RSIAT gốc $74.76/65.77$.
\item Hai bản dùng chung backbone (đặc trưng giống hệt); mọi khác biệt đến từ thống kê lớp và cách phân loại.
\end{itemize}

% =====================================================================================================================
\begin{thebibliography}{99}
\bibitem{slca} G. Zhang, L. Wang, G. Kang, L. Chen, Y. Wei. SLCA: Slow Learner with Classifier Alignment for Continual
Learning on a Pre-trained Model. \emph{ICCV}, 2023.
\bibitem{ssiat} Y. Tan, Q. Zhou, X. Xiang, K. Wang, Y. Wu, Y. Li. Semantically-Shifted Incremental Adapter-Tuning is A
Continual ViTransformer. \emph{CVPR}, 2024.
\bibitem{sdc} L. Yu, B. Twardowski, X. Liu, L. Herranz, K. Wang, Y. Cheng, S. Jui, J. van de Weijer. Semantic Drift
Compensation for Class-Incremental Learning. \emph{CVPR}, 2020.
\bibitem{adc} D. Goswami, A. Soutif-Cormerais, Y. Liu, S. Kamath, B. Twardowski, J. van de Weijer. Resurrecting Old
Classes with New Data for Exemplar-Free Continual Learning. \emph{CVPR}, 2024.
\bibitem{ldc} A. Gomez-Villa, D. Goswami, K. Wang, A. D. Bagdanov, B. Twardowski, J. van de Weijer. Exemplar-free
Continual Representation Learning via Learnable Drift Compensation. \emph{ECCV}, 2024.
\bibitem{cosface} H. Wang, Y. Wang, Z. Zhou, X. Ji, D. Gong, J. Zhou, Z. Li, W. Liu. CosFace: Large Margin Cosine Loss
for Deep Face Recognition. \emph{CVPR}, 2018.
\bibitem{esl} T. Hastie, R. Tibshirani, J. Friedman. \emph{The Elements of Statistical Learning}, 2nd ed., Springer, 2009.
\bibitem{rahimi} A. Rahimi, B. Recht. Random Features for Large-Scale Kernel Machines. \emph{NeurIPS}, 2007.
\bibitem{procrustes} P. H. Schönemann. A Generalized Solution of the Orthogonal Procrustes Problem. \emph{Psychometrika},
1966.
\bibitem{murphy} K. P. Murphy. \emph{Machine Learning: A Probabilistic Perspective}. MIT Press, 2012.
\end{thebibliography}

\end{document}
